\documentclass{article}
\title{Economia e Mercados Financeiros}
\begin{document}
\maketitle

\section{Juros e Inflacao}
A taxa real de juros aproxima a diferenca entre a taxa nominal e a
inflacao esperada. Formalmente, a relacao de Fisher e dada por:
\[ (1 + i) = (1 + r)(1 + \pi) \]
onde $i$ e a taxa nominal, $r$ a taxa real e $\pi$ a inflacao.

\section{Valor Presente}
O valor presente de um fluxo futuro $C$ descontado a uma taxa $r$ por
$n$ periodos e $PV = \frac{C}{(1 + r)^n}$. Para varios fluxos:

\[ PV = \sum_{t=1}^{n} \frac{C_t}{(1 + r)^t} \]

\subsection{Pontos principais}
\begin{itemize}
\item Juros compostos crescem de forma exponencial.
\item Inflacao corroi o poder de compra ao longo do tempo.
\item O risco exige um premio sobre a taxa livre de risco.
\end{itemize}
\end{document}
